\section{总结与延伸}

\subsection{Post-Training 流程总览}

整个 Post-Training 是一个迭代循环，由三个等重的环节构成：

\begin{enumerate}
    \item \textbf{数据构建}（约 1/3 时间）：设计数据混合比例，确保准确性、多样性和复杂度，通过 seed data + LLM 生成 + 过滤的流水线批量生产高质量训练数据
    \item \textbf{模型训练}（约 1/3 时间）：SFT（推荐 LoRA）$\rightarrow$ Preference Alignment（推荐 DPO）$\rightarrow$ 可选的 Model Merging
    \item \textbf{评估迭代}（约 1/3 时间）：自动化 benchmark + LLM-as-Judge + 人类评估三管齐下，评估结果反馈到下一轮数据和训练的改进
\end{enumerate}

\subsection{核心要点回顾}

\begin{importantbox}{本讲核心 Takeaways}
\begin{itemize}
    \item Post-Training 将 base model 转化为有用的 assistant，核心步骤为 SFT + Preference Alignment
    \item 数据质量 $>$ 数据数量，三个维度：Accuracy, Diversity, Complexity
    \item LoRA 是大多数场景下的最佳 SFT 技术，DPO 是最实用的对齐算法
    \item Model Merging 是零成本组合多模型优势的利器
    \item 人类偏好与 automated benchmark 相关性低，必须组合使用多种评估方法
    \item Test-Time Compute Scaling 用推理时间换输出质量，是当前最重要的趋势
\end{itemize}
\end{importantbox}

\subsection{方法选择的工程决策表}

\begin{table}[H]
\centering
\begin{tabular}{p{0.2\textwidth} p{0.24\textwidth} p{0.24\textwidth}}
\toprule
目标场景 & 优先方法 & 关键判断标准 \\
\midrule
先快速验证需求 & In-context learning / RAG & 不要过早进入 fine-tuning，先看 base model 是否已足够好 \\
需要稳定输出格式或语气 & SFT（通常配 LoRA） & 样本质量是否足够高，输出风格是否可被明确示范 \\
需要让回答更符合人类偏好 & DPO / Preference Alignment & 是否有 chosen vs rejected 的偏好数据，评估是否覆盖真实用户偏好 \\
需要组合多个模型专长 & Model Merging & 各模型能力是否互补，以及合并后是否有退化风险 \\
推理阶段还有预算 & Test-Time Compute Scaling & 延迟和成本是否允许用更多采样换更高质量 \\
\bottomrule
\end{tabular}
\caption{Post-Training 主要方法的选择顺序}
\end{table}

\subsection{给实践者的落地顺序}

如果把这一讲转成工程上的执行建议，最稳妥的顺序通常是：

\begin{enumerate}
    \item 先用最小成本的 prompting / RAG 确认问题边界
    \item 再用高质量小数据集做 SFT，解决风格和格式问题
    \item 若要进一步提升可用性，再引入偏好数据做 DPO
    \item 最后根据成本、延迟与效果，选择 model merging 或 test-time scaling 作为增强手段
\end{enumerate}

\subsection{拓展阅读}

\begin{itemize}
    \item \textbf{LLM Engineer's Handbook}——讲者撰写的系统性参考书
    \item \textbf{Open Perfect Blend} (Hugging Face)——通用 SFT 数据集混合方案
    \item \textbf{mergekit}——开源模型合并工具
    \item \textbf{TRL} (Hugging Face)——最全面的 Post-Training 工具库
    \item Rafailov et al., \textit{Direct Preference Optimization} (NeurIPS 2023)——DPO 原始论文
    \item Hu et al., \textit{LoRA: Low-Rank Adaptation of Large Language Models} (ICLR 2022)——LoRA 原始论文
    \item Snell et al., \textit{Scaling LLM Test-Time Compute} (2024)——Test-Time Compute 的理论基础
    \item Chatbot Arena (\url{https://arena.lmsys.org})——基于人类偏好的模型排名
\end{itemize}

%% --- Fin del contenido --- %%

\end{document}
